\chapter{Theory and Concept}
\label{ch:theory}

This chapter sets out the concepts the benchmark is built on. It defines the
vulnerability classes in scope, the families of detectors that the benchmark
scores, the label-noise problem that motivates the validation methodology, the
notion of semantics-preserving robustness, and the evaluation measures used to
report results. It closes with a review of related work across datasets, label
noise, detectors, and robustness studies.

%----------------------------------------------------------

\section{Sink-Shaped Vulnerability Classes}

A vulnerability class in the Common Weakness Enumeration (CWE) catalog describes
a recurring kind of security flaw. This work targets seven \emph{sink-shaped}
classes, where a sink is a function whose unsafe use is the vulnerability and
whose safe use is not. SQL injection (CWE-89) occurs when untrusted input
reaches a database query sink without parameterization. Cross-site scripting
(CWE-79) reaches an HTML rendering sink, OS command injection (CWE-78) reaches a
process-execution sink, code injection (CWE-94) reaches \texttt{eval} or
\texttt{exec}, path traversal (CWE-22) reaches a filesystem sink, server-side
request forgery (CWE-918) reaches an outbound HTTP sink, and insecure
deserialization (CWE-502) reaches a sink such as \texttt{yaml.load} or
\texttt{pickle.loads}.

The sink-shaped property matters because it gives each positive an anchor: the
fix for the vulnerability modifies the sink call. Structural classes, such as
missing authorization or cross-site request forgery, encode the \emph{absence}
of a check and require application-level context to judge. They have no single
line that a diff can anchor on, so they are excluded from this benchmark on
methodological grounds rather than as a statement about their importance.

\section{Detector Families}

The benchmark scores detectors from four tiers, each resting on a different
principle.

\begin{itemize}
  \item \textbf{Pattern-matching SAST.} Bandit \cite{bandit} and Semgrep
    \cite{semgrep} scan a single file for syntactic signatures of unsafe
    constructs, in the same lexical family as general Python linters such as
    Pylint \cite{pylint}. They are fast and need no build, but they reason about
    the surface form of code rather than its data flow.
  \item \textbf{Whole-project dataflow.} CodeQL \cite{codeql} and Snyk Code
    track data from a source to a sink across an entire application. They are
    more precise on complete programs but need a built database with
    application entry points.
  \item \textbf{Large language models.} LLMs reason over code in natural
    language and are competitive on C and C++ benchmarks
    \cite{khare_llm,steenhoek}, but their verdicts can shift under small input
    changes, and code they generate carries security bugs of its own
    \cite{pearce2022}.
  \item \textbf{Learned classifiers.} A model such as GraphCodeBERT
    \cite{graphcodebert,codebert} is fine-tuned on labeled vulnerability data
    and learns features from both token sequence and data-flow structure,
    following a line of deep vulnerability detectors that begins with token-level
    \cite{vuldeepecker} and representation-learning \cite{russell2018} models.
\end{itemize}

\section{The Label-Noise Problem}

Detector training and evaluation depend on labels that say which code is
vulnerable. Most labels are inherited from CVE fix commits. A fix commit
touches several files for refactoring, tests, or cleanup that do not relate to
the vulnerability, and each co-changed file inherits the CWE label of
the advisory. Croft et al. \cite{croft} audited BigVul and reported about 25\%
noise; the ReVeal authors \cite{reveal} showed that models learn co-changed
surface patterns rather than the vulnerability itself. CleanVul \cite{cleanvul}
found 40\% to 75\% of labels inaccurate across common datasets, and
PrimeVul \cite{primevul} showed that a model scoring 68.3\% F1 on BigVul drops
to 3.1\% on a deduplicated, tightened-label version of the same task. The
lesson is that a benchmark built on raw label inheritance measures a detector's
agreement with label noise as much as its detection ability. The validation
methodology in Chapter~\ref{ch:impl} addresses this with a sink-anchored,
line-level criterion.

\section{Semantics-Preserving Robustness}

A reported detection score is measured on unmodified code. The robustness
question asks whether that score survives a rewrite of the input that changes
its surface form but not its runtime behavior. A semantics-preserving
transformation is one that leaves the program's observable behavior unchanged,
for example renaming a local variable or splitting a string literal. If a
detector's verdict changes under such a transformation, the detector relied on a
surface shortcut rather than on the program's meaning. SecLLMHolmes
\cite{secllmholmes} showed that such edits flip LLM verdicts, and the broader
literature on code-model robustness \cite{yefet,bielik,henkel} establishes that
learned models are sensitive to these edits. The mutators in this work differ
by targeting one interpretable shortcut each, so a per-mutator accuracy drop is
diagnostic rather than a single opaque attack score.

\section{Evaluation Measures}

Detection quality is reported with five bounded, complementary measures rather
than a single composite, so a reader can see where a detector is strong. Macro
averaged F1 \cite{sokolova2009} summarizes per-class precision and recall across
the seven classes, which matters because class sizes are unequal. The Matthews
correlation coefficient (MCC) \cite{chicco2020mcc} scores the binary
vulnerable-versus-safe decision and holds up when the classes are balanced.
Severity-Weighted Recall (SWR) weights recall by CVSS severity so that missing a
critical flaw costs more than missing a minor one. Hierarchical CWE Macro
Accuracy (HCMA) gives Wu-Palmer \cite{wupalmer1994} partial credit for
close-sibling confusions on a CWE subtree. Weighted Cohen's $\kappa$
\cite{cohen1968kappa} reports chance-corrected agreement. Robustness is reported
as the drop in macro-F1 between clean and mutated code, computed over the same
positives so the subtraction is like-for-like.

\section{Related Work}

\subsection{Datasets}

Source-code vulnerability benchmarks are C and C++ dominant. BigVul
\cite{bigvul}, CVEfixes \cite{cvefixes}, CrossVul \cite{crossvul}, and the
deduplicated DiverseVul \cite{diversevul} contribute few per-class Python
samples. Juliet \cite{juliet} is synthetic and does not transfer to natural code
\cite{reveal}; Devign and ReVeal \cite{devign,reveal} are C only. Within Python,
\texttt{vudenc} \cite{vudenc} is the closest labeled corpus but ships no paired
fixed versions, which constrains the label validation it admits.

\subsection{Label Noise and Validation}

Beyond the audits noted above, D2A \cite{d2a} uses differential analysis for
labels but operates at program rather than sink granularity. The diff-based
filter in this work is a sink-anchored, line-level label-validity criterion,
and to our knowledge it is the first such evidence-based criterion in this
literature.

\subsection{Detectors and Robustness}

The most relevant prior comparison is CASTLE \cite{castle}, which pits static
analyzers against LLMs for C and C++ but uses CVE-inherited labels and adds no
robustness axis. RealVuln \cite{realvuln} targets Python repositories with
manual labels but reports detection only. IRIS \cite{iris} is an LLM and CodeQL
hybrid whose CWE-Bench-Java targets the same sink-shaped family. LineVul
\cite{linevul} and MoreFixes \cite{morefixes} extend the dataset and modeling
line. SecLLMHolmes \cite{secllmholmes} established that semantics-preserving
edits flip detector verdicts, which motivates the robustness axis here.

\subsection{Positioning}

Each ingredient has precedent. Large vulnerability datasets, the identification
of CVE label noise, the SAST-versus-LLM comparison, and code-model robustness
studies all exist on their own. This work joins them for Python: seven
sink-shaped Top-25 CWE classes, an automated sink-anchored label-validity
criterion, and a per-mutator robustness breakdown reported next to detection
accuracy in one evaluation. That last step surfaces which detectors reason about
code and which match its surface.
